\par\noindent\textbf{(1)} 第二London方程为 \(\nabla\times\boldsymbol J_s=-(n_se^2/m)\boldsymbol B\)，静态Ampère方程 \(\nabla\times\boldsymbol B=\mu_0\boldsymbol J_s\)。取旋度并利用 \(\nabla\cdot\boldsymbol B=0\)，得 \(\nabla^2\boldsymbol B=\boldsymbol B/\lambda_L^2\)，其中 \(\lambda_L^2=m/(\mu_0n_se^2)\)。
\par\noindent\textbf{(2)} 由于场只依赖 \(x\)，满足 \(B_z''=B_z/\lambda_L^2\)。半空间不发散的解为 \(B_z(x)=B_0e^{-x/\lambda_L}\)。由 \((\nabla\times\boldsymbol B)_y=-B_z'(x)\) 得 \(J_y(x)=B_0e^{-x/\lambda_L}/(\mu_0\lambda_L)\)，与屏蔽外场的面电流方向一致。
\par\noindent\textbf{(3)} 每单位面积磁场储能 \(U_B/A=\int_0^\infty B_z(x)^2/(2\mu_0)\,dx=B_0^2\lambda_L/(4\mu_0)\)。总表面电流 \(K_y=\int_0^\infty J_y\,dx=B_0/\mu_0\)，与边界磁场跃变相符；London动能另外贡献能量，不能把纯磁场积分冒充完整超导自由能。
\par\noindent\textbf{(4)} 当 \(\lambda_L\to0\)，场趋于表面薄层而内部 \(B_z(x>0)\to0\)，即理想Meissner屏蔽。实际第二类超导体超过下临界场 \(H_{c1}\) 可进入涡旋混合态，非局域响应、有限厚度及临界电流也会使简单指数解失效。
